\section{Proposta Metodológica}
    \label{sec:proposta-metodologica}

    Os pesquisadores selecionados para a pesquisa de campo foram os responsáveis pelo desenvolvimento do projeto \textit{SmartEnergy}, que consiste em um sistema de monitoramento de consumo energia elétrica, com algoritmos de aprendizado de máquina para detecção de padrões e anomalias.

    Os pesquisadores foram entrevistados para determinar as principais características que o software deveria possuir. A partir das respostas obtidas, pode ser determinados os requisitos necessários para o desenvolvimento do software. Com base nesses requisitos, o software proposto pode ser desenvolvido. Os pesquisadores envolvidos levantaram questões como \textit{Qual a melhor forma de visualizar os dados?}, \textit{Como integrar novos algoritmos?}, \textit{Como exportar os resultados?}, \textit{Como obter estatísticas numéricas dos dados?}, \textit{Qual a melhor forma de realizar a importação dos dados?}, entre outras. Todas as questões foram avaliadas e as respostas aceitas foram utilizadas para a definição dos requisitos do software.

    A partir dos requisitos levantados, foi realizado a pesquisa bibliográfica para identificar projetos semelhantes e conhecer as soluções implementadas. A partir da pesquisa notamos uma variedade de softwares semelhantes, em linguagens como \textbf{Java}, \textbf{R}, \textbf{Octave} e \textbf{MATLAB}. Linguagens como as citadas, são tipicamente utilizadas por pesquisadores da área, pelo fato de já possuírem uma gama de algoritmos destinados à análise de dados já implementados.

    Visto que o principal motivo para o uso das linguagens citadas é a existência de bibliotecas de algoritmos famosos já implementados, o Python foi escolhido para o desenvolvimento deste trabalho. Como já apontado anteriormente, o Python possui uma vasta biblioteca, com uma comunidade de usuário ativa e crescente. Além disso, por possuir uma sintaxe simples, ela é ideal para que os pesquisadores, e futuros usuários, possam criar módulos capazes de estender as funcionalidades do projeto.

    Por fim, a avaliação do software foi realizada com os pesquisadores. Eles propuseram a utilização e estudo de uma base de dados para avaliar a capacidade do sistema, de modo que seja satisfatório o problema proposto por ela.
